\section{Evaluation}
\label{sec:eval}
A document can appear in the catalogue yet be missing from the index. A
search can reach the right collection yet miss the passage that answers the
question. We test these failures separately from answer quality, which also
depends on how much text the answering agent can read.

\subsection{Datasets and measures}\label{sec:eval-setup}
The answer comparison uses ten developer-authored questions with annotated
reference papers in five arXiv collections. The routing and evidence-pack
experiments use 72 in-domain questions across ten collections (138 documents
in the frozen run); out-of-domain probes are analysed separately. The
verifier diagnostics use a separate industrial corpus of 42 documents in
five collections. The saved questions omit the gold document identifiers.

For routing, a hit means that the selector chose the labelled collection.
For evidence packs, a hit means that the annotated target string occurs in
the returned text. This second measure does not verify a claim and can count
a match in a different document from the annotated source. Hit@5 measures
whether a gold document appears among the first five results. Answer scores
and verdicts are defined in Section~\ref{sec:e10}. The paired sign and
McNemar tests are two-sided and exploratory. Repeated runs use the same
questions; they do not add independent observations.

Table~\ref{tab:ablation-configs} lists the six configurations used to test
how evidence packs are assembled.
\begin{table}[!ht]
\centering
\small
\setlength{\tabcolsep}{3pt}
\renewcommand{\arraystretch}{1.16}
\caption{Search scope and evidence construction for the six E4 configurations.
Limits are in characters, not model tokens.}
\label{tab:ablation-configs}
\begin{tabular}{@{}l>{\raggedright\arraybackslash}p{0.23\columnwidth}>{\raggedright\arraybackslash}p{\dimexpr0.77\columnwidth-31pt\relax}@{}}
\toprule
Arm & Search scope & Candidate selection and evidence\\
\midrule
A0 & One collection, selected by name and description & Retrieve vector hits. The reference-mode gate reads up to six distinct documents, with a nominal 40{,}000-character window per document.\\
A0f & Same route as A0 & Use the same candidates with an inline gate. Allow up to 250{,}000 characters per candidate and in the gate pack.\\
A0h & Same route as A0 & Keep the A0f gate and add a second vector search. Place up to eight chunk windows first, then add gate-retained text up to a final limit of 60{,}000 characters.\\
PD & All ten collections & Keep the six highest-scoring distinct collection/path pairs and use the A0f inline gate.\\
A1b & All ten collections & Merge vector hits and use the A0 reference-mode gate, including document reads.\\
A2 & All ten collections & Merge the same hits as A1b and concatenate their snippets. Do not read documents or run a gate.\\
\bottomrule
\end{tabular}
\end{table}

Routed vector search requests up to 30 candidates at a cosine threshold of
0.35 and retries at 0.30 only if it returns nothing. The unrouted arms search
each collection at 0.35. Some reference lookups could use the gold domain,
but it did not determine the route-card choice (Section~\ref{sec:threats}).

\subsection{Do changes in storage reach the index? (E1)}
\label{sec:e1}
\label{sec:e2}
After explicit indexing and reindexing, all 328 stored arXiv parts had
content-derived descriptions, and search probes returned identifiers for
all 100 papers. This establishes identifier presence, not ranking within
each collection. The controlled checks passed for document creation, moves,
deletion and forced reindexing. In an earlier disposable-collection test,
a rename left 18 vector entries instead of nine. Forced reindexing restored
the expected count. In a separate, later retest, an old label in search
results was traced to the retrieval cache; the index metadata was correct.
Task-status and update checks failed; merge and split were not tested. The
historical tag audit also recorded orphan tags and a cross-collection
registry anomaly. These checks distinguish what is stored, what is indexed
and what the search interface returns.
